\documentclass[11pt,a4paper]{coverletter}

\senderblock
  {Radheem Bin Razi}
  {Distributed Systems \& Cloud-Native Engineer}
  {Ilmenau, Germany \textbar\ \href{mailto:sheikh.radheem@gmail.com}{sheikh.radheem@gmail.com} \textbar\ \href{https://radheem.github.io/my-cv}{portfolio}}

\begin{document}

%% ===== ENGLISH =====
\recipient{payabl}{Hiring Team}{}

\opening{Dear Hiring Team,}

\vspace{8pt}\noindent\textbf{1. Why payabl?}\par\vspace{3pt}

Paying systems live or die by consistency and latency. I enjoy designing backend services that keep transactions flowing smoothly while quietly handling the messy edge cases that users never see. That focus on resilient, high-throughput payment infrastructure is exactly why I am applying for the Backend Engineer role at payabl.

\vspace{8pt}\noindent\textbf{2. Why me?}\par\vspace{3pt}

My background centers on building distributed backends that scale under real load. At a stealth tax platform, I replaced a heavy sidecar architecture with native NATS JetStream for messaging and direct gRPC calls, which cut deployment complexity and improved message routing reliability. I also designed an anti-corruption adaptor layer to safely integrate multiple bot-protected third-party vendors, ensuring schema mismatches never reached the core service.

I have spent the last few years writing production Python services and event-driven pipelines that feed high-throughput data stores. I containerized ingestion flows and built Kafka-based metrics pub/sub systems that fan out telemetry to multiple databases, giving teams clear visibility into system health. I am comfortable working with CI/CD pipelines, writing unit tests that catch regressions early, and collaborating with DevOps teams to keep deployments predictable. I also maintain strict code quality standards by reviewing peer pull requests and documenting architectural decisions. These experiences align directly with your need for scalable microservices, third-party API integration, and reliable production operations.

\vspace{8pt}\noindent\textbf{3. Why now?}\par\vspace{3pt}

I am looking to join a team that treats payment infrastructure as a craft, and payabl’s focus on expanding high-performance backend capabilities matches my next career step perfectly. I am available to start full-time immediately and am prepared to relocate to Cyprus. I hold a Blue Card-ready status, which means your team can issue a standard employment contract without any visa sponsorship overhead. I look forward to discussing how my backend experience can support your growing engineering squad.

\closing{Sincerely,}{Radheem Bin Razi}

\clearpage
\selectlanguage{ngerman}

%% ===== DEUTSCH =====
\recipient{payabl}{Personalabteilung}{}

\opening{Sehr geehrtes Hiring-Team,}

\vspace{8pt}\noindent\textbf{1. Warum payabl?}\par\vspace{3pt}

Zahlungssysteme leben oder sterben an Konsistenz und Latenz. Ich gestalte gerne Backend-Services, die Transaktionen reibungslos ablaufen lassen und dabei stillschweigend die komplexen Edge Cases abfangen, die Nutzer nie zu sehen bekommen. Genau dieser Fokus auf eine resiliente, hochdurchsatzstarke Zahlungsinfrastruktur ist der Grund, warum ich mich für die Position als Backend Engineer bei payabl bewerbe.

\vspace{8pt}\noindent\textbf{2. Warum ich?}\par\vspace{3pt}

Meine Expertise liegt im Aufbau verteilter Backends, die auch unter realer Last skalieren. Bei einer Stealth-Tax-Plattform habe ich eine schwere Sidecar-Architektur durch native NATS JetStream für Messaging und direkte gRPC-Aufrufe ersetzt, was die Komplexität der Bereitstellung reduzierte und die Zuverlässigkeit der Nachrichtenweiterleitung verbesserte. Zudem entwarf ich eine Anti-Corruption-Adapter-Schicht, um mehrere bot-geschützte Drittanbieter sicher zu integrieren und sicherzustellen, dass Schema-Inkonsistenzen niemals den Core-Service erreichen.

In den letzten Jahren habe ich Production-Python-Services und ereignisgesteuerte Pipelines entwickelt, die hochdurchsatzstarke Datenspeicher versorgen. Ich containerisierte Ingestion-Flows und baute Kafka-basierte Metrics-Pub/Sub-Systeme auf, die Telemetrie-Daten auf mehrere Datenbanken verteilen und den Teams so eine klare Übersicht über die Systemgesundheit ermöglichen. Ich verfüge über umfassende Erfahrung in der Arbeit mit CI/CD-Pipelines, beim Schreiben von Unit Tests, die Regressionen frühzeitig erkennen, sowie bei der Zusammenarbeit mit DevOps-Teams, um Bereitstellungen stabil und vorhersehbar zu halten. Zudem halte ich strikte Code-Qualitätsstandards ein, indem ich Peer-Pull Requests überprüfe und architektonische Entscheidungen dokumentiere. Diese Erfahrungen entsprechen direkt Ihrem Bedarf an skalierbaren Microservices, Drittanbieter-API-Integrationen und einem zuverlässigen Production-Betrieb.

\vspace{8pt}\noindent\textbf{3. Warum jetzt?}\par\vspace{3pt}

Ich suche ein Team, das Zahlungsinfrastruktur als Handwerk mit höchster Sorgfalt gestaltet, und der Fokus von payabl auf den Ausbau hochperformanter Backend-Kapazitäten passt perfekt zu meinem nächsten Karriereschritt. Ich stehe ab sofort zur Verfügung und kann das Vollzeit-Engagement aufnehmen; ich bin bereit, nach Cyprus zu ziehen. Ich verfüge über einen Blue Card-ready Status, was bedeutet, dass Ihr Team einen regulären Arbeitsvertrag ohne zusätzlichen Visums-Sponsorship-Aufwand ausstellen kann. Ich freue mich darauf, in einem Gespräch zu besprechen, wie meine Backend-Erfahrung Ihr wachsendes Engineering-Team unterstützen kann.

\closing{Mit freundlichen Grüßen,}{Radheem Bin Razi}

\end{document}
